\small
\begin{enumerate}[label=\textbf{\arabic*.}, leftmargin=0pt, itemsep=1.2em]
\item Для чего служит директива :- dynamic bigger/2. в начале программы?
\begin{itemize}[label=$\square$, leftmargin=1.5cm, itemsep=0.2em, topsep=0.3em]
\item Включает трассировку предиката bigger/2
\item Объявляет, что клаузы bigger/2 можно добавлять и удалять во время выполнения
\item Запрещает изменять факты bigger/2 после загрузки файла
\item Делает предикат bigger/2 видимым из других файлов
\end{itemize}
{\small\textit{Правильный ответ: Объявляет, что клаузы bigger/2 можно добавлять и удалять во время выполнения.}}

\item Чем assertz/1 отличается от asserta/1?
\begin{itemize}[label=$\square$, leftmargin=1.5cm, itemsep=0.2em, topsep=0.3em]
\item assertz удаляет клаузу, а asserta добавляет
\item assertz добавляет клаузу в конец предиката, asserta — в начало
\item assertz работает только с фактами, asserta — только с правилами
\item Ничем, это синонимы
\end{itemize}
{\small\textit{Правильный ответ: assertz добавляет клаузу в конец предиката, asserta — в начало.}}

\item Что делает вызов retractall(bigger(\_, \_))?
\begin{itemize}[label=$\square$, leftmargin=1.5cm, itemsep=0.2em, topsep=0.3em]
\item Удаляет только первый факт bigger/2
\item Удаляет все клаузы bigger/2, при этом сам предикат остаётся определённым
\item Выводит все факты bigger/2 на экран
\item Завершает работу интерпретатора
\end{itemize}
{\small\textit{Правильный ответ: Удаляет все клаузы bigger/2, при этом сам предикат остаётся определённым.}}

\item Чем retract/1 отличается от retractall/1?
\begin{itemize}[label=$\square$, leftmargin=1.5cm, itemsep=0.2em, topsep=0.3em]
\item retract удаляет первую подходящую клаузу и при возврате — следующие, retractall удаляет все подходящие за один вызов
\item retract удаляет предикат целиком, retractall — только факты
\item retractall можно применять только к статическим предикатам
\item Ничем
\end{itemize}
{\small\textit{Правильный ответ: retract удаляет первую подходящую клаузу и при возврате — следующие, retractall удаляет все подходящие за один вызов.}}

\item Что вернёт read\_line\_to\_string(user\_input, L) при достижении конца ввода (Ctrl+D)?
\begin{itemize}[label=$\square$, leftmargin=1.5cm, itemsep=0.2em, topsep=0.3em]
\item Пустую строку ""
\item Атом halt
\item Ошибку синтаксиса
\item Атом end\_of\_file
\end{itemize}
{\small\textit{Правильный ответ: Атом end\_of\_file.}}

\item Каков результат split\_string("яблоко > слива", " ", "", P)?
\begin{itemize}[label=$\square$, leftmargin=1.5cm, itemsep=0.2em, topsep=0.3em]
\item P = "яблоко > слива"
\item P = [яблоко, >, слива]
\item P = ["яблоко", ">", "слива"]
\item false.
\end{itemize}
{\small\textit{Правильный ответ: P = ["яблоко", ">", "слива"].}}

\item Для чего в программе используется normalize\_space(string(Clean), Line)?
\begin{itemize}[label=$\square$, leftmargin=1.5cm, itemsep=0.2em, topsep=0.3em]
\item Переводит строку в нижний регистр
\item Убирает лишние пробелы в начале, конце и между словами
\item Разбивает строку на слова
\item Преобразует строку в число
\end{itemize}
{\small\textit{Правильный ответ: Убирает лишние пробелы в начале, конце и между словами.}}

\item Зачем после format("ввод: ") вызывают flush\_output?
\begin{itemize}[label=$\square$, leftmargin=1.5cm, itemsep=0.2em, topsep=0.3em]
\item Чтобы очистить экран терминала
\item Чтобы закрыть поток вывода
\item Чтобы приглашение без перевода строки немедленно появилось на экране до чтения ввода
\item Чтобы удалить введённые ранее данные
\end{itemize}
{\small\textit{Правильный ответ: Чтобы приглашение без перевода строки немедленно появилось на экране до чтения ввода.}}

\item Что делает вызов prompt(Old, '') перед чтением строк из user\_input?
\begin{itemize}[label=$\square$, leftmargin=1.5cm, itemsep=0.2em, topsep=0.3em]
\item Запоминает прежнее приглашение в Old и заменяет его пустым, чтобы не выводилось «|: »
\item Запрещает вводить пустые строки
\item Очищает буфер ввода
\item Загружает файл с пустым именем
\end{itemize}
{\small\textit{Правильный ответ: Запоминает прежнее приглашение в Old и заменяет его пустым, чтобы не выводилось «|: ».}}

\item Зачем у правила greater/2 две клаузы?
\begin{itemize}[label=$\square$, leftmargin=1.5cm, itemsep=0.2em, topsep=0.3em]
\item Вторая клауза нужна для обработки ошибок
\item Первая — базовый случай (прямой факт bigger), вторая — рекурсивный шаг по цепочке отношений
\item Первая работает с числами, вторая — с атомами
\item Две клаузы нужны только для ускорения поиска
\end{itemize}
{\small\textit{Правильный ответ: Первая — базовый случай (прямой факт bigger), вторая — рекурсивный шаг по цепочке отношений.}}

\item Что произошло бы с greater/2 при цикле в фактах (a больше b, b больше a)?
\begin{itemize}[label=$\square$, leftmargin=1.5cm, itemsep=0.2em, topsep=0.3em]
\item Запрос вернул бы false
\item Интерпретатор отказался бы загружать факты
\item Рекурсия не имела бы конца и привела бы к переполнению стека
\item Результат зависел бы только от порядка ввода
\end{itemize}
{\small\textit{Правильный ответ: Рекурсия не имела бы конца и привела бы к переполнению стека.}}

\item Какую роль выполняет клауза add\_relation(X, Y) :- greater(Y, X), !, ...?
\begin{itemize}[label=$\square$, leftmargin=1.5cm, itemsep=0.2em, topsep=0.3em]
\item Добавляет новый факт в базу
\item Обнаруживает противоречие: Y уже больше X, и отклоняет новый факт
\item Удаляет старые факты
\item Сортирует объекты по размеру
\end{itemize}
{\small\textit{Правильный ответ: Обнаруживает противоречие: Y уже больше X, и отклоняет новый факт.}}

\item Как читается правило largest(X) :- object(X), \textbackslash{}+ bigger(\_, X).?
\begin{itemize}[label=$\square$, leftmargin=1.5cm, itemsep=0.2em, topsep=0.3em]
\item X больше любого объекта
\item X — объект, для которого не найдётся ни одного объекта, больший X
\item X — объект, который больше самого себя
\item X — любой объект, имеющий хотя бы один факт bigger
\end{itemize}
{\small\textit{Правильный ответ: X — объект, для которого не найдётся ни одного объекта, больший X.}}

\item Что означает оператор \textbackslash{}+ в Прологе?
\begin{itemize}[label=$\square$, leftmargin=1.5cm, itemsep=0.2em, topsep=0.3em]
\item Арифметическое сложение
\item Логическое отрицание в смысле математической логики
\item Отрицание как неуспех: цель \textbackslash{}+ G успешна, если доказать G не удаётся
\item Объединение двух списков
\end{itemize}
{\small\textit{Правильный ответ: Отрицание как неуспех: цель \textbackslash{}+ G успешна, если доказать G не удаётся.}}

\item Что делает предикат findall(X, largest(X), L)?
\begin{itemize}[label=$\square$, leftmargin=1.5cm, itemsep=0.2em, topsep=0.3em]
\item Находит первое решение и останавливается
\item Удаляет все решения largest/1 из базы
\item Собирает в список L значения X по всем решениям largest(X)
\item Проверяет, что largest/1 определён
\end{itemize}
{\small\textit{Правильный ответ: Собирает в список L значения X по всем решениям largest(X).}}

\item Почему для примера из задания (яблоко, слива, авокадо, папайя, дуриан, вишня) нет единого самого большого объекта?
\begin{itemize}[label=$\square$, leftmargin=1.5cm, itemsep=0.2em, topsep=0.3em]
\item Потому что вводились не все объекты
\item Потому что дуриан меньше папайи
\item Яблоко и дуриан не связаны цепочкой отношений, и ни один из них не доказуемо больше другого
\item Потому что в Прологе нельзя сравнивать более пяти объектов
\end{itemize}
{\small\textit{Правильный ответ: Яблоко и дуриан не связаны цепочкой отношений, и ни один из них не доказуемо больше другого.}}

\item Какой самый маленький объект в примере с фруктами из задания?
\begin{itemize}[label=$\square$, leftmargin=1.5cm, itemsep=0.2em, topsep=0.3em]
\item слива
\item авокадо
\item вишня
\item яблоко
\end{itemize}
{\small\textit{Правильный ответ: вишня.}}

\item Для чего служит вызов aggregate\_all(count, bigger(\_, \_), M)?
\begin{itemize}[label=$\square$, leftmargin=1.5cm, itemsep=0.2em, topsep=0.3em]
\item Суммирует значения аргументов bigger/2
\item Подсчитывает число решений цели, то есть число фактов bigger/2
\item Удаляет дубликаты фактов
\item Сортирует факты по возрастанию
\end{itemize}
{\small\textit{Правильный ответ: Подсчитывает число решений цели, то есть число фактов bigger/2.}}

\item Какой арности предикат день\_рождения(иванова, лена, 12, июнь, 1988)?
\begin{itemize}[label=$\square$, leftmargin=1.5cm, itemsep=0.2em, topsep=0.3em]
\item 3
\item 4
\item 5
\item 6
\end{itemize}
{\small\textit{Правильный ответ: 5.}}

\item Какие люди вернёт запрос \texttt{?-} день\_рождения(Ф, И, \_, \_, 1988).?
\begin{itemize}[label=$\square$, leftmargin=1.5cm, itemsep=0.2em, topsep=0.3em]
\item Только Лену Иванову
\item Только Васю Сидоренко
\item Лену Иванову и Машу Сидоренко
\item Всех троих друзей
\end{itemize}
{\small\textit{Правильный ответ: Лену Иванову и Машу Сидоренко.}}

\item Что означает символ \_ в запросе \texttt{?-} день\_рождения(Ф, И, \_, май, \_).?
\begin{itemize}[label=$\square$, leftmargin=1.5cm, itemsep=0.2em, topsep=0.3em]
\item Пустое значение, которое обязательно нужно подставить
\item Имя предиката
\item Анонимная переменная: значение на этой позиции не важно и не запоминается
\item Знак вычитания
\end{itemize}
{\small\textit{Правильный ответ: Анонимная переменная: значение на этой позиции не важно и не запоминается.}}

\item Сколько решений даст запрос \texttt{?-} хобби(\_, маша, Х).?
\begin{itemize}[label=$\square$, leftmargin=1.5cm, itemsep=0.2em, topsep=0.3em]
\item 0
\item 1
\item 2
\item 3
\end{itemize}
{\small\textit{Правильный ответ: 2.}}

\item Почему запрос \texttt{?-} хобби(Ф, И, кино), хобби(Ф, И, спорт). находит именно Васю?
\begin{itemize}[label=$\square$, leftmargin=1.5cm, itemsep=0.2em, topsep=0.3em]
\item Потому что Вася родился раньше всех
\item Одни и те же переменные Ф и И в обеих целях требуют, чтобы оба хобби принадлежали одному человеку
\item Потому что запятая означает «или»
\item Потому что Вася — единственный друг в базе
\end{itemize}
{\small\textit{Правильный ответ: Одни и те же переменные Ф и И в обеих целях требуют, чтобы оба хобби принадлежали одному человеку.}}

\item Какое имя в SWI-Prolog является переменной?
\begin{itemize}[label=$\square$, leftmargin=1.5cm, itemsep=0.2em, topsep=0.3em]
\item иванова
\item 'Лена'
\item Ф
\item маша
\end{itemize}
{\small\textit{Правильный ответ: Ф.}}

\item Как называется метод решения логической задачи в программе sportsmen.pl?
\begin{itemize}[label=$\square$, leftmargin=1.5cm, itemsep=0.2em, topsep=0.3em]
\item Динамическое программирование
\item Метод наименьших квадратов
\item Жадный алгоритм
\item Образовать и проверить (generate and test)
\end{itemize}
{\small\textit{Правильный ответ: Образовать и проверить (generate and test).}}

\item Сколько различных результатов даёт запрос \texttt{?-} findall(P, permutation([a, b, c], P), L), length(L, N).?
\begin{itemize}[label=$\square$, leftmargin=1.5cm, itemsep=0.2em, topsep=0.3em]
\item 3
\item 6
\item 9
\item 27
\end{itemize}
{\small\textit{Правильный ответ: 6.}}

\item Зачем для видов спорта и для погоды используют permutation/2, а не независимый выбор значений?
\begin{itemize}[label=$\square$, leftmargin=1.5cm, itemsep=0.2em, topsep=0.3em]
\item Для ускорения работы интерпретатора
\item Перестановка гарантирует, что значения признака у троих спортсменов различны и все использованы
\item Чтобы получить значения в алфавитном порядке
\item Это требуется синтаксисом SWI-Prolog
\end{itemize}
{\small\textit{Правильный ответ: Перестановка гарантирует, что значения признака у троих спортсменов различны и все использованы.}}

\item Какое условие задачи выражает цель \textbackslash{}+ member(человек(\_, боксёр, солнце), Люди)?
\begin{itemize}[label=$\square$, leftmargin=1.5cm, itemsep=0.2em, topsep=0.3em]
\item Боксёр любит только солнечную погоду
\item Тот, кто любит солнечную погоду, не боксёр
\item Среди спортсменов нет людей
\item Все любят солнечную погоду
\end{itemize}
{\small\textit{Правильный ответ: Тот, кто любит солнечную погоду, не боксёр.}}

\item Кем является Тонни согласно найденному решению?
\begin{itemize}[label=$\square$, leftmargin=1.5cm, itemsep=0.2em, topsep=0.3em]
\item Футболистом, любит снег
\item Пловцом, любит солнце
\item Футболистом, любит дождь
\item Боксёром, любит дождь
\end{itemize}
{\small\textit{Правильный ответ: Боксёром, любит дождь.}}

\item Что подтверждает результат «Найдено решений: 1» для логической задачи?
\begin{itemize}[label=$\square$, leftmargin=1.5cm, itemsep=0.2em, topsep=0.3em]
\item Что программа завершилась без ошибок и это всё
\item Что условия задачи определяют ответ однозначно
\item Что в задаче нет ограничений
\item Что программа проверила только один вариант
\end{itemize}
{\small\textit{Правильный ответ: Что условия задачи определяют ответ однозначно.}}
\end{enumerate}
\normalsize
